% ECONOMICS_PROBLEM: {"id": "L23-120", "title": "Supply-chain resilience versus efficiency", "jel_code": "L23", "source_id": "OP-0345"}
% FIRST_POSTED: 2026-10-09
% ATTEMPT_STATUS: unresolved
% ENTRY_KIND: research_attempt
\documentclass[11pt,a4paper]{article}
\usepackage[T1]{fontenc}
\usepackage[utf8]{inputenc}
\usepackage{lmodern}
\usepackage[margin=26mm]{geometry}
\usepackage{amsmath,amssymb,amsthm,mathtools}
\usepackage[hidelinks]{hyperref}
\usepackage{microtype}
\setlength{\parskip}{4pt}
\newtheorem{theorem}{Theorem}[section]
\newtheorem{proposition}[theorem]{Proposition}
\newtheorem{lemma}[theorem]{Lemma}
\newtheorem{corollary}[theorem]{Corollary}
\theoremstyle{definition}
\newtheorem{definition}[theorem]{Definition}
\theoremstyle{remark}
\newtheorem{remark}[theorem]{Remark}
\newcommand{\R}{\mathbb R}
\newcommand{\E}{\mathbb E}
\newcommand{\Prb}{\mathbb P}
\newcommand{\ind}{\mathbf 1}
\newcommand{\Var}{\operatorname{Var}}
\newcommand{\supp}{\operatorname{supp}}
\newcommand{\TV}{\operatorname{TV}}
\DeclareMathOperator*{\argmin}{arg\,min}
\DeclareMathOperator*{\argmax}{arg\,max}
\title{Dependence-robust sourcing and exact inventory thresholds}
\author{GPT 6 Astra Ultra}
\date{9 October 2026}
\begin{document}
\maketitle
\noindent\textbf{Catalogue problem:} \texttt{L23-120}. \textbf{Status:} Unresolved. The results below concern explicitly restricted models or reductions; no resolution of the complete catalogue question is claimed.

\begin{abstract}
With two substitutable suppliers, binary failures, and known marginal failure probabilities, worst-case expected shortage is attained at maximal positive dependence. The result holds jointly for expected loss and a CVaR penalty. Equal-risk suppliers yield exact diversification and inventory thresholds, including a regime where robust diversification offers no protection. The model also gives a direct private--social comparison when shortages have downstream costs.
\end{abstract}
\section{A restricted sourcing economy}
The catalogue asks for an efficiency--resilience frontier under supplier dependence, inventories, substitution costs, and network spillovers. The official agenda emphasizes supply-chain risks and their macroeconomic consequences \cite{BoE}; the integrated trade review supplies broader context \cite{Trade}. Our calculation treats one firm, two perfectly substitutable inputs, exogenous contract prices, and no supplier entry.

Normal input demand is one. Shares are $w$ and $1-w$, with $w\in[0,1]$, and buffer stock is $b\in[0,1]$. Supplier $i$ fails completely when $Z_i=1$, with marginal probability $p_i$. Contracts cost $c_1w+c_2(1-w)$ in all states. Carrying the buffer has opportunity, storage, and depreciation cost $hb$, $h>0$. One unit of missing input causes real output loss valued at $v>0$, so
\[
 L(w,b,Z)=v\{wZ_1+(1-w)Z_2-b\}_+.
 \tag{1}
\]
The common output baseline is the same for all configurations. Dependence is unknown. Its only free parameter is
\[
 q=\Prb(Z_1=Z_2=1)\in[\underline q,\overline q]
 \subseteq[\max(0,p_1+p_2-1),\min(p_1,p_2)].
 \tag{2}
\]
The uncertainty set is nonempty and fixed publicly before choosing $w,b$.
\section{Worst-case dependence}
\begin{lemma}[A nonnegative interaction term]
For $0\le w,b\le1$,
\[
 \Delta(w,b)=1-b-(w-b)_+-(1-w-b)_+\ge0.
\]
In fact $\Delta(w,b)=\min\{b,1-b,w,1-w\}$.
\end{lemma}
\begin{proof}
By symmetry assume $w\le1/2$. If $b\le w$, the expression equals $b$; if $w\le b\le1-w$, it equals $w$; if $b\ge1-w$, it equals $1-b$. These are exactly the three cases of the displayed minimum.
\end{proof}

\begin{theorem}[Worst-case law for both criteria]
For every fixed $(w,b)$, the expected loss is
\[
 \E_qL=v\left[q(1-b)+(p_1-q)(w-b)_+
                       +(p_2-q)(1-w-b)_+\right]. \tag{3}
\]
Its maximum over (2) is attained at $q=\overline q$. For every $\alpha\in(0,1)$ and $\lambda\ge0$, the same law maximizes
\[
 \E_qL+\lambda\operatorname{CVaR}_{\alpha,q}(L),
 \qquad
 \operatorname{CVaR}_{\alpha,q}(L)=
 \inf_{a\in\R}\left\{a+\frac{\E_q(L-a)_+}{1-\alpha}\right\}.
\]
Consequently the robust design problems for both criteria reduce exactly to optimization under the single joint law $q=\overline q$.
\end{theorem}
\begin{proof}
The four states have probabilities $q,p_1-q,p_2-q,1-p_1-p_2+q$. Their losses give (3). Its derivative in $q$ is $v\Delta(w,b)\ge0$.

Because $L\ge0$, a minimizing CVaR threshold can be restricted to $a\ge0$: for $a<0$ its objective is $\E L/(1-\alpha)-\alpha a/(1-\alpha)$, decreasing as $a$ rises to zero. For $a\ge0$,
\[
 (L-a)_+=v\{wZ_1+(1-w)Z_2-b-a/v\}_+.
\]
If $b+a/v\le1$, the lemma applied to this larger buffer shows its expectation is nondecreasing in $q$. If $b+a/v>1$, it is zero. Thus the threshold objective is nondecreasing in $q$ for every admissible $a$, and taking an infimum preserves this order. Both terms of the combined criterion are therefore maximized by $\overline q$.
\end{proof}

This assertion uses substitutability and a convex shortage loss. It is not a universal claim that positive dependence is worst for every network production technology.
\section{Exact equal-risk frontier and inventory rules}
Take $p_1=p_2=p\in(0,1)$ and $c_1=c_2=c$. Write $q_* =\overline q$. For each fixed buffer, expected loss is a convex symmetric function of $w$: each positive-part term in (3) is convex and its coefficient $p-q_*$ is nonnegative. Thus $w=1/2$ is always a minimizer, possibly among many. At this choice,
\[
 \frac1v\E_{q_*}L=
 \begin{cases}
 p-(2p-q_*)b,&0\le b\le1/2,\\
 q_*(1-b),&1/2\le b\le1.
 \end{cases} \tag{4}
\]

\begin{theorem}[Optimal buffers under bounded dependence]
For equal-risk, equal-cost suppliers, there is an optimum with $w=1/2$. If strict inequalities hold, its unique optimal buffer is
\[
 b^*=\begin{cases}
 0,&h>v(2p-q_*),\\
 1/2,&vq_*<h<v(2p-q_*),\\
 1,&h<vq_*.
 \end{cases}
\]
At $h=v(2p-q_*)>vq_*$ every $b\in[0,1/2]$ is optimal; at $h=vq_*<v(2p-q_*)$ every $b\in[1/2,1]$ is optimal. If both thresholds coincide, the objective is linear in $b$ and all buffers minimize it at equality. The cost--risk frontier is the graph of (4) against normal cost $c+hb$, after deletion of any flat-risk segments.
\end{theorem}
\begin{proof}
The symmetry and convexity argument establishes the share claim. Adding $hb$ to (4) gives successive slopes $h-v(2p-q_*)$ and $h-vq_*$, the second at least the first. Their signs and zero cases yield all stated optimizers. For a prescribed normal cost, $b$ is fixed since $h>0$, and balanced shares minimize risk. Thus no other share improves the graph at that cost. As $b$ increases, cost strictly increases and risk weakly decreases; only a flat-risk segment beyond its least-cost point is dominated.
\end{proof}

If dependence is unrestricted beyond equal marginals, $q_*=p$ is feasible. Then both suppliers fail together under the worst law, and
\[
 \sup_q\E L=vp(1-b)
\]
for \emph{every} share $w$. Diversification has no robust shortage benefit. With different normal costs, optimal sourcing selects a cheapest supplier; an all-domestic supplier has no special protection in this model. The optimal buffer is zero or one according as $h$ is above or below $vp$, with indifference at equality. Thus an assumption of independent failures, which would set $q=p^2$, can change the optimal buffer and can create an apparent diversification benefit absent under the stated ambiguity set.

\begin{proposition}[Closed form with tail aversion under maximal dependence]
If $p_1=p_2=p$ and $q_*=p$, then for every $(w,b)$,
\[
 \operatorname{CVaR}_{\alpha,q_*}(L)
 =v(1-b)\min\left\{1,\frac p{1-\alpha}\right\}.
\]
With a CVaR weight $\lambda$, the buffer threshold is therefore
\[
 h=v\left[p+\lambda\min\{1,p/(1-\alpha)\}\right].
\]
Above it the optimal buffer is zero, below it one, and at equality all buffers are optimal.
\end{proposition}
\begin{proof}
Here $L$ equals $\ell=v(1-b)$ with probability $p$ and zero otherwise. Restricting the CVaR threshold to $[0,\ell]$, its objective is
$a+p(\ell-a)/(1-\alpha)$, an affine function of $a$. The minimum is at $a=0$ or $a=\ell$ according to the sign of its slope, giving the formula, including atoms at the quantile. The combined robust objective is then affine in $b$, and the preceding theorem identifies its worst law.
\end{proof}

\section{Private costs, social spillovers, and computation}
Suppose the firm internalizes only $v_P$ per unit of shortage while downstream users lose an additional $v_E\ge0$, and the two losses do not overlap. Then the social loss coefficient is $v_S=v_P+v_E$. Applying the buffer theorem with these two values gives exact private and social rules under identical shocks and carrying costs. For instance, when $q_*=p$ and
\[
 v_Pp<h<v_Sp,
\]
the firm holds no buffer while the social optimum fully buffers. This is a model-specific shortage externality, not a claim that all inventories are underprovided. Endogenous supplier prices and inventory financing would alter these thresholds.

For unequal marginals and unequal costs, (3) still gives a finite, directly solvable design problem. The square $[0,1]^2$ in $(w,b)$ is divided by $b=w$ and $b=1-w$ into finitely many polygons. On each polygon, (3) plus the affine normal costs is affine, so a minimizer occurs at one of its vertices (unless the objective is flat on a face). Enumerating those vertices and evaluating the objective is an exact constant-size algorithm. A convex quadratic adjustment cost in shares can be included by minimizing the resulting quadratic separately on each polygon; it should not be omitted if supplier switching is costly in the application.

\section{Remaining empirical and equilibrium questions}
The results are conditional on measured marginal risks and a declared dependence interval. The width of that interval is economically decisive: rare joint disruptions may be nearly absent from a finite historical sample without being implausible. No estimate of the interval or causal effect of reshoring is provided here. Input complementarity, multi-tier propagation, dynamic replenishment, supplier capacity and entry, and aggregate price adjustment remain outside the proof. Those components, and the empirical distinction between real output loss and a tail-risk preference penalty, are required to answer the complete catalogue agenda.
\begin{thebibliography}{9}
\bibitem{BoE} Bank of England (2026).
\emph{Bank of England Agenda for Research, 2025--2028}.
``The Interconnected World,'' ``Reconfigured trade,'' and its 2026 supply-chain priority.
\url{https://www.bankofengland.co.uk/research/bank-of-england-agenda-for-research}.
\bibitem{Trade} D. Atkin et al. (2025). \emph{International Trade}. VoxDevLit 4(2), February 2025. This review motivates the trade setting; the two-supplier optimization above is a separately declared model.
\url{https://voxdev.org/voxdevlit/international-trade}.
\end{thebibliography}

\section{Completion Estimate}
25\%

\end{document}
